\documentclass[a4paper]{article}
\usepackage[pdftex]{graphicx}
\usepackage[utf8]{inputenc}
\usepackage{enumerate}
\usepackage{colortbl}
\usepackage{href-ul}
\hypersetup{
	colorlinks=true,
	linkcolor=blue,
	urlcolor=blue}
\usepackage{siunitx}
\sisetup{locale=DE} 
\usepackage{icomma}
\usepackage{amssymb}
\usepackage{geometry}
\geometry{a4paper, top=15mm, left=15mm, right=15mm, bottom=15mm,
	headsep=10mm, footskip=12mm}

\begin{document}
{\bf Kinetische Energie}

\begin{enumerate}[1.]
{\bf \item Gegeben sind folgende Körper mit charakteristischen Massen und Geschwindigkeiten.  Ordne sie der Tabelle zu.}:\\
\fbox{Fußball}, \fbox{Sternschnuppe}, \fbox{Auto}, \fbox{ICE}, \fbox{Tischtennisball}, \fbox{Kampfjet},\fbox{Regentropfen},\fbox{Sprinter}

\renewcommand{\arraystretch}{2}
\begin{tabular}{|>{\columncolor[gray]{.8}}p{3 cm}|p{3 cm}|p{3cm}|p{3 cm}|}
	\hline
	\rowcolor[gray]{.8}	Körper & Masse &Geschwindigkeit & Kinetische Energie \\
	\hline
	  &$ \SI{0,05}{\gram} $&$ \SI{5}{\meter\per\second} $&\\
	\hline
	&$ \SI{4,7}{\gram} $&$ \SI{10}{\meter\per\second} $& \\
	\hline
	&$ \SI{75}{\kilo\gram} $&$ \SI{10}{\meter\per\second} $&\\
	\hline
	&$ \SI{0,1}{\gram} $&$ \SI{30}{\kilo\meter\per\second} $& \\
	\hline
	&$ \SI{0,42}{\kilo\gram} $&$ \SI{100}{\kilo\meter\per\hour} $& \\
	\hline
	&$ \SI{1}{\tonne} $&$ \SI{130}{\kilo\meter\per\hour} $&  \\	
	\hline
	&$ \SI{2}{\tonne} $&$ \SI{1200}{\kilo\meter\per\hour} $&\\
	\hline
	&$ \SI{300}{\tonne} $&$ \SI{250}{\kilo\meter\per\hour} $&\\
	\hline
\end{tabular}



{\bf \item Ein Wagen der Masse $m=\SI{200}{\kilogram}$ einer Achterbahn durchfährt einen Looping.} Er startet dazu mit der Geschwindigkeit $v_0 = 0$ im Punkt A in einer Höhe von $h=\SI{30}{\meter}$ (Siehe Skizze, nicht maßstabsgerecht). 

\begin{enumerate}[a)]
\item Berechne die Geschwindigkeit des Wagens im Punkt B und D (ohne Berücksichtigung der Reibung)
\item Wie ändern sich diese Geschwindigkeiten, wenn der Wagen mehr als $\SI{200}{\kilogram}$ Masse hätte  (ohne Berücksichtigung der Reibung)?
\item Der Wagen soll hinter D auf einer Strecke von $s=\SI{25}{\meter}$ zum Stillstand gebracht werden. Wie groß ist die erforderliche Bremskraft?
\end{enumerate}

\includegraphics[width=9 cm]{achterbahn115.pdf}


{\bf \item Aufgabe\\}
\begin{enumerate}[a)]
\item Berechne die Geschwindigkeit, mit der ein Turmspringer beim Sprung vom $\SI{10}{\meter}$ - Turm
auf das Wasser auftrifft.
\item Aus welcher Höhe ist jemand gesprungen, der eine Geschwindigkeit von $v=\SI{10}{\meter\per\second}$ kurz vor dem Eintauchen ins Wasser erreicht hat?
\end{enumerate} 
\end{enumerate} 

\fbox{
	\begin{minipage}{0.5\textwidth}
		Zur Lösung bitte \href{https://www.okuyakl.de/phys/p8ekiL411/ll411.pdf}{hier klicken} oder den QR-Code scannen.\\
	Weitere Arbeitsblätter gibt es unter 
	
	\href{https://www.okuyakl.de}{www.okuyakl.de}
	\end{minipage}
	\hfill
	\begin{minipage}{0.4\textwidth}
		\includegraphics[width=1.5 cm]{../../viecher/zwe03}
		\includegraphics[width=3 cm]{qr411}
		\includegraphics[width=2 cm]{../../viecher/afanticon1}
		
	\end{minipage}}
\end{document}
